\section{The contour and the trial measure}\label{sec:contour-and-measure}

This appendix specifies the polygon $\Gamma$ and the trial measure $\mu$ to which the energy
criterion (\cref{prop:energy-criterion}) is applied in \cref{prop:energy-certificate}, and records
the properties of $\mu$ that the criterion requires. The data are stored in the file \file{trial.json} of the original archive (SHA-256
\texttt{a4da00b3a51e1d41\ldots c499c}); all coordinates are rational with common denominator
$7\,054\,320\,000\,000$ and all masses are integers divided by $10^{12}$.

\subsection{The polygon}\label{subsec:polygon}

The upper half $\Gamma^+$ of $\Gamma$ is the polygonal path through the vertices of
\cref{tab:vertices}, followed by the vertical ray $v_7+i[0,\infty)$; $\Gamma=\Gamma^+\cup\overline{\Gamma^+}$.
The verifier checks that all vertices have negative real part, that $v_0$ is real and that the
ordinates increase strictly, so $\Gamma$ is a simple curve in $\{\operatorname{Re}z<0\}$ meeting
$\R$ only at $v_0=-r_0$.

\begin{table}[ht]
\centering
\small
\begin{tabular}{@{}cll@{}}
\toprule
vertex & real part & imaginary part\\
\midrule
$v_0$ & $-664907/5000000=-0.1329814$ & $0$\\
$v_1$ & $-61937/400000=-0.1548425$ & $1/25$\\
$v_2$ & $-11119/62500=-0.177904$ & $9/100$\\
$v_3$ & $-2046703/10000000=-0.2046703$ & $4/25$\\
$v_4$ & $-583527/2500000=-0.2334108$ & $1/4$\\
$v_5$ & $-170547/625000=-0.2728752$ & $2/5$\\
$v_6$ & $-170547/625000$ & $3/4$\\
$v_7$ & $-170547/625000$ & $5/4$\\
\bottomrule
\end{tabular}
\caption{Vertices of the upper half of the scaled contour $\Gamma$.}
\label{tab:vertices}
\end{table}

\subsection{The trial measure}\label{subsec:trial-measure}

The trial measure $\mu$ is a sum of $238$ uniform measures on segments:
\begin{itemize}
  \item $34$ consecutive segments of equal length partitioning $[-r_0,0]$, carrying total mass
  $0.2659627999\ldots$, with density between $2-6\cdot10^{-10}$ and $2-3\cdot10^{-11}$;
  \item $102$ segments of length between $0.0038$ and $0.0040$ on the edges of $\Gamma^+$ up to
  height $0.377$, with densities between $0.028$ and $1.92$ and total mass
  $0.3982686000\ldots$;
  \item the $102$ complex conjugate segments, with the same masses.
\end{itemize}
The verifier checks exactly: every mass is positive, the total mass is $17/16=\lambda$, the
measure is invariant under conjugation, every segment lies on an edge of $\Gamma$ or on
$[v_0,0]$, and every endpoint $w$ satisfies $\abs{\operatorname{Re}w}+\abs{\operatorname{Im}w}\le1$,
so $\operatorname{supp}\mu\subset\{\abs w\le1\}$. Thus $\mu$ satisfies the hypotheses of
\cref{prop:energy-criterion}.

The trial measure was produced by a floating-point optimizer (a discretized quadratic program for
the constrained equilibrium of \cref{rem:energy-meaning}) and then rounded to rationals; the
certificate does not depend on how it was found. In the AI-internal audit, half of the $3639$
cell endpoints of the certificate have a value of $V$ within about $1.5\cdot10^{-4}$ of the
minimum. These endpoints are concentrated on the part of $\Gamma$ that carries $\mu$, of height at
most $0.377$; above it $V$ increases, up to about $4.12$ on the finite part of the contour. So $V$
is nearly constant on the support of $\mu$ in $\Gamma$, as it would be for an equilibrium measure;
this suggests that $\mu$ is close to optimal for this contour, but we have not proved it and do
not use it. The audit also finds $V<\theta$ on all but about $6\cdot10^{-5}$ of the interval
$[-r_0,0]$, so the term $4J\approx-0.239$ matters.
